\documentclass[11pt,a4paper]{article}
\usepackage[portuguese]{babel}

\usepackage{array, tabularx, multirow, xcolor, colortbl, makecell}


\usepackage[utf8]{inputenc}
\usepackage[T1]{fontenc}
\usepackage[portuguese]{babel}
\usepackage{array}
\usepackage{geometry}
\usepackage{longtable}
\usepackage{booktabs}
\usepackage{multirow}
\geometry{margin=2cm}

\definecolor{headerpink}{RGB}{220,80,120}
\definecolor{lightpink}{RGB}{255,230,240}



\RequirePackage{graphicx}
\RequirePackage{tikz}

\providecommand{\PastaLogosAlfornelos}{assets/files/}

\newcommand{\CabecalhoAlfornelos}{%
  \begingroup
  \setlength{\parindent}{0pt}
  % --- Dimensões do bloco (mantém a proporção do cabeçalho original) ---
  \def\LarguraCab{\textwidth}
  \newlength\AlturaCab
  \setlength{\AlturaCab}{0.1525\LarguraCab}
  %
  \noindent


  \begin{tikzpicture}
   \useasboundingbox (0,0) rectangle (\LarguraCab,-1.42\AlturaCab);

     % --- Logótipo esquerdo ---
    \node[anchor=north west, inner sep=0] at (0.0308\LarguraCab, -0.1709\AlturaCab)
      {\includegraphics[width=0.1805\LarguraCab]{\PastaLogosAlfornelos logo_republica_portuguesa_completo.png}};

    % --- Texto das escolas (usa a fonte corrente do documento) ---
    \node[anchor=north, align=center, text width=0.99\LarguraCab, inner sep=0]
    at (0.5\LarguraCab, -0.2730\AlturaCab)
   {\fontsize{9}{11}\selectfont\bfseries
    Agrupamento de Escolas de Alfornelos\\
    Código 170161 \\
    A Par e Passo Rumo ao Sucesso\\
};

 % --- Logótipo agrupamento ---
    \node[anchor=north west, inner sep=0] at (0.8525\LarguraCab, -0.1519\AlturaCab)
      {\includegraphics[width=0.1971\LarguraCab]{\PastaLogosAlfornelos logo_agrupamento_completo.png}};

      % --- Linha inferior ---
    \draw[line width=0.9pt, black!70]
      (0.1342\LarguraCab, -0.9658\AlturaCab) -- (0.8668\LarguraCab, -0.9658\AlturaCab);


    %--- Texto inferior ---
    \node[anchor=north, align=center, text width=0.99\LarguraCab, inner sep=0]
    at (0.5\LarguraCab, -1\AlturaCab)
   {\fontsize{9}{11}\selectfont\bfseries
   Sede: Escola Básica de Alfornelos\\
   344515

};
  \end{tikzpicture}
  \endgroup
}


\begin{document}
\CabecalhoAlfornelos

\newcolumntype{L}{>{\small\bfseries}p{2.5cm}}
\newcolumntype{Y}{>{\raggedright\arraybackslash}X}

\renewcommand{\arraystretch}{1.2}

\begin{tabularx}{\textwidth}{|>{\bfseries\raggedright\arraybackslash}p{3.5cm}|X|}
\hline
Escola & Escola Básica de Alfornelos \\
\hline
Ano letivo & 2026/2027 \\
\hline
Serviço & Departamento de Matemática e Informática \\
\hline
Grupo disciplinar & 550 - Informática \\
\hline
Assunto & Planificação Anual de TIC \\
\hline
Professor & Filipe Onofre\\
\hline
\end{tabularx}

\renewcommand{\arraystretch}{1.3}

\rowcolors{2}{lightpink}{white}
\begin{longtable}{|p{3cm}|p{4cm}|p{3.5cm}|p{3.5cm}|p{3.5cm}|p{2cm}|}
\hline
\rowcolor{headerpink}
\textbf{Organizador / Domínio} &
\textbf{AE: Conhecimentos, Capacidades e Atitudes} &
\textbf{Conteúdos} &
\textbf{Operacionalização} &
\textbf{Avaliação} &
\textbf{Tempos (min)} \\
\hline
\endfirsthead

\hline
\rowcolor{headerpink}
\textbf{Organizador / Domínio} &
\textbf{AE: Conhecimentos, Capacidades e Atitudes} &
\textbf{Conteúdos} &
\textbf{Operacionalização} &
\textbf{Avaliação} &
\textbf{Tempos (min)} \\
\hline
\endhead

Segurança, responsabilidade e respeito em ambientes digitais. &
O aluno adota uma atitude crítica e responsável no uso de tecnologias e serviços digitais. &
Princípios da sociedade digital; práticas seguras; direitos de autor; regras de ergonomia. &
Debates em grupo; atividades interativas; apresentação de trabalhos. &
Critérios: conhecimentos adquiridos, aplicação, participação, interesse, criatividade. &
50 \\
\hline

% Adiciona mais linhas conforme necessário
% Exemplo:
% Outro domínio &
% Descrição das atitudes &
% Conteúdos &
% Estratégias &
% Avaliação &
% Tempo \\

\end{longtable}



\end{document}